\documentclass[12pt]{article}
\usepackage[a4paper,margin=1in]{geometry}
\usepackage[utf8]{inputenc}
\usepackage[T1]{fontenc}
\usepackage{amsmath,amssymb}
\usepackage{textcomp}
\usepackage{parskip}
\begin{document}
\begin{center}
{\Large \textbf{2000 STEP I}}
\end{center}

\subsection*{Page 1}
\begin{flushleft}
STEP I, 2000 2\par
\bigskip\textbf{Section A: Pure Mathematics}\par
\bigskip\textbf{Question 1.} To nine decimal places, log 2 = 0.301029996 and log 3 = 0.477121255.\par
10 10\par
\textbf{(i)} Calculate log 5 and log 6 to three decimal places. By taking logs, or otherwise,\par
10 10\par
show that\par
47 100 47\par
\[5 \times 10 < 3 < 6 \times 10 .\]
100\par
Hence write down the first digit of 3 .\par
1000 10 000 100 000\par
\textbf{(ii)} Find the first digit of each of the following numbers: 2 ; 2 ; and 2 .\par
\bigskip\textbf{Question 2.} Show that the coefficient of x -12 in the expansion of\par
5 6\par
1 1\par
x - x -\par
\[µ x^2 ¶ µ x ¶\]
is -15, and calculate the coefficient of x .\par
4 38\par
Hence, or otherwise, calculate the coefficients of x and x in the expansion of\par
2 11 4 2 5\par
(x - 1) (x + x + 1) .\par
\bigskip\textbf{Question 3.} For any number x, the largest integer less than or equal to x is denoted by [x]. For example,\par
\[[3.7] = 3 and [4] = 4.\]
\[Sketch the graph of y = [x] for 0 6 x < 5 and evaluate\]
[x] dx.\par
\[Z^0\]
\[Sketch the graph of y = [ e x ] for 0 6 x < ln n, where n is an integer, and show that\]
ln n\par
x\par
\[[ e ] dx = n ln n - ln(n!).\]
\[Z^0\]
\end{flushleft}
\subsection*{Page 2}
\begin{flushleft}
STEP I, 2000 3\par
x 1\par
\bigskip\textbf{Question 4.} (i) Show that, for 0 6 x 6 1, the largest value of is .\par
\[(x^2 + 1)4 16\]
\textbf{(ii)} Find constants A, B, C and D such that, for all x,\par
5 3 6\par
1 d Ax + Bx + Cx Dx\par
\[= + .\]
\[(x^2 + 1)4 dx µ (x^2 + 1)3 ¶ (x^2 + 1)4\]
\textbf{(iii)} Hence, or otherwise, prove that\par
11 1 11 1\par
6 dx 6 + .\par
\[24 Z^0 (x^2 + 1)4 24 16\]
5 Arthur and Bertha stand at a point O on an inclined plane. The steepest line in the plane\par
through O makes an angle θ with the horizontal. Arthur walks uphill at a steady pace in a\par
straight line which makes an angle α with the steepest line. Bertha walks uphill at the same\par
speed in a straight line which makes an angle β with the steepest line (and is on the same\par
side of the steepest line as Arthur). Show that, when Arthur has walked a distance d, the\par
distance between Arthur and Bertha is 2d| sin (α - β)|. Show also that, if α 6= β, the line\par
joining Arthur and Bertha makes an angle φ with the vertical, where\par
\[cos φ = sin θ sin (α + β).\]
\bigskip\textbf{Question 6.} Show that\par
2 2\par
\[x - y + x + 3y - 2 = (x - y + 2)(x + y - 1)\]
and hence, or otherwise, indicate by means of a sketch the region of the x-y plane for which\par
2 2\par
\[x - y + x + 3y > 2.\]
Sketch also the region of the x-y plane for which\par
2 2\par
\[x - 4y + 3x - 2y < -2.\]
Give the coordinates of a point for which both inequalities are satisfied or explain why no such\par
point exists.\par
\end{flushleft}
\subsection*{Page 3}
\begin{flushleft}
STEP I, 2000 4\par
7 Let\par
x\par
\[f (x) = ax - ,\]
\[1 + x^2\]
\[where a is a constant. Show that, if a > 9/8, then f 0 (x) > 0 for all x.\]
\bigskip\textbf{Question 8.} Show that\par
1 0 1\par
x x x\par
\[| x e | dx = - x e dx + x e dx\]
\[Z-1 Z-1 Z^0\]
and hence evaluate the integral.\par
Evaluate the following integrals:\par
3 2\par
\textbf{(i)} | x - 2x - x + 2 | dx ;\par
\[Z^0\]
\[\pi\]
\textbf{(ii)} | sin x + cos x | dx .\par
\[Z-\pi\]
\end{flushleft}
\subsection*{Page 4}
\begin{flushleft}
STEP I, 2000 5\par
\bigskip\textbf{Section B: Mechanics}\par
\bigskip\textbf{Question 9.} A child is playing with a toy cannon on the floor of a long railway carriage. The carriage is\par
moving horizontally in a northerly direction with acceleration a. The child points the cannon\par
southward at an angle θ to the horizontal and fires a toy shell which leaves the cannon at\par
speed V . Find, in terms of a and g, the value of tan 2θ for which the cannon has maximum\par
range (in the carriage).\par
If a is small compared with g, show that the value of θ which gives the maximum range is\par
approximately\par
\[\pi a\]
+ ,\par
4 2g\par
2 2\par
V V a\par
and show that the maximum range is approximately + .\par
\[g g^2\]
\bigskip\textbf{Question 10.} Three particles P1, P2 and P3 of masses m1, m2 and m3 respectively lie at rest in a straight\par
\[line on a smooth horizontal table. P^1 is projected with speed v towards P^2 and brought to\]
rest by the collision. After P2 collides with P3, the latter moves forward with speed v. The\par
coefficients of restitution in the first and second collisions are e and e , respectively. Show that\par
\[m^2 + m^3 - m^1\]
\[e = .\]
\[m^1\]
\[Show that 2m1 > m^2 + m^3 > m^1 for such collisions to be possible.\]
If m1, m3 and v are fixed, find, in terms of m1, m3 and v, the largest and smallest possible\par
values for the final energy of the system.\par
\bigskip\textbf{Question 11.} A rod AB of length 0.81 m and mass 5 kg is in equilibrium with the end A on a rough floor\par
and the end B against a very rough vertical wall. The rod is in a vertical plane perpendicular\par
\textasciicircum{}\circ\par
to the wall and is inclined at 45 to the horizontal. The centre of gravity of the rod is at G,\par
where AG = 0.21 m. The coefficient of friction between the rod and the floor is 0.2, and the\par
coefficient of friction between the rod and the wall is 1.0. Show that the friction cannot be\par
limiting at both A and B.\par
A mass of 5 kg is attached to the rod at the point P such that now the friction is limiting at\par
both A and B. Determine the length of AP .\par
\end{flushleft}
\subsection*{Page 5}
\begin{flushleft}
STEP I, 2000 6\par
\bigskip\textbf{Section C: Probability and Statistics}\par
\bigskip\textbf{Question 12.} I have k different keys on my key ring. When I come home at night I try one key after another\par
until I find the key that fits my front door. What is the probability that I find the correct key\par
in exactly n attempts in each of the following three cases?\par
\textbf{(i)} At each attempt, I choose a key that I have not tried before but otherwise each choice\par
is equally likely.\par
\textbf{(ii)} At each attempt, I choose a key from all my keys and each of the k choices is equally\par
likely.\par
\textbf{(iii)} At the first attempt, I choose from all my keys and each of the k choices is equally\par
likely. Thereafter, I choose from the keys that I did not try the previous time but\par
otherwise each choice is equally likely.\par
13 Every person carries two genes which can each be either of type A or of type B. It is known\par
that 81\% of the population are AA (i.e. both genes are of type A), 18\% are AB (i.e. there is\par
one gene of type A and one of type B) and 1\% are BB. A child inherits one gene from each\par
of its parents. If one parent is AA, the child inherits a gene of type A from that parent; if the\par
parent is BB, the child inherits a gene of type B from that parent; if the parent is AB, the\par
inherited gene is equally likely to be A or B.\par
\textbf{(i)} Given that two AB parents have four children, show that the probability that two of\par
\[them are AA and two of them are BB is 3/128.\]
\textbf{(ii)} My mother is AB and I am AA. Find the probability that my father is AB.\par
\bigskip\textbf{Question 14.} The random variable X is uniformly distributed on the interval [-1, 1]. Find E(X 2 ) and\par
Var (X ).\par
A second random variable Y , independent of X, is also uniformly distributed on [-1, 1], and\par
2 2 2\par
\[Z = Y - X. Find E(Z ) and show that Var (Z ) = 7 Var (X ).\]
\end{flushleft}
\end{document}
